\section{从 InstructGPT 到 ChatGPT}

\subsection{InstructGPT：定义管线}

InstructGPT 是第一个系统性地将前述所有技术组合在一起的模型，它定义了现代 LLM 后训练的标准管线。

\begin{figure}[H]
\centering
\includegraphics[width=0.95\textwidth]{slides-images/slide-078.jpg}
\caption{InstructGPT 的完整管线：SFT $\rightarrow$ 奖励模型 $\rightarrow$ PPO\protect\footnotemark}
\end{figure}
\footnotetext{Slides 第79页。}

InstructGPT 的关键数据：
\begin{itemize}
    \item 约 \textbf{30,000 种不同任务}的指令微调数据
    \item 使用人类标注者收集偏好排序
    \item 通过 PPO 进行 RLHF 优化
\end{itemize}

\begin{figure}[H]
\centering
\includegraphics[width=0.95\textwidth]{slides-images/slide-079.jpg}
\caption{InstructGPT 效果：GPT-3 对``向6岁小孩解释登月''的回答 vs InstructGPT 的回答\protect\footnotemark}
\end{figure}
\footnotetext{Slides 第80页。}

效果对比非常直观：GPT-3 面对``向6岁小孩解释登月''的问题，会继续生成更多问题（因为它只是在补全文本）；而 InstructGPT 则直接给出适合 6 岁儿童理解的简明解释。

\subsection{ChatGPT：对话优化}

ChatGPT 在 InstructGPT 的基础上，进一步针对\textbf{对话交互}进行了优化：

\begin{figure}[H]
\centering
\includegraphics[width=0.95\textwidth]{slides-images/slide-085.jpg}
\caption{ChatGPT 的对话能力：理解用户意图，生成有趣且相关的回答\protect\footnotemark}
\end{figure}
\footnotetext{Slides 第86页。}

核心技术不变（SFT + RLHF），关键差异在于数据——更多的对话格式数据、更注重多轮交互的连贯性。

\subsection{本章小结}

InstructGPT 和 ChatGPT 证明了后训练管线的有效性：通过 SFT 建立指令遵循能力，通过 RLHF/DPO 进一步对齐人类偏好。这个管线已经成为现代 LLM 开发的标准流程。

%% ========== Part 7: 挑战与展望 ==========

